% ==========================================================================
%  Simulado P2 — Prova + Gabarito Comentado
%  Macroeconomia I (PPGEco-UFES, 2026)
%  Modelo: P1 2026 (escolha única — julgar afirmativas V/F)
%  Cobre: Consumo (cap.8), Investimento (cap.9),
%         Inflação e Política Monetária (cap.11), Política Fiscal (cap.12)
%  Fontes: P2 2021, P2 2024 (chave oficial), Lista 3 2026 (chave oficial),
%          slides das Aulas 10–13
% ==========================================================================
\documentclass[12pt, a4paper]{article}

\usepackage[utf8]{inputenc}
\usepackage[T1]{fontenc}
\usepackage{lmodern}
\usepackage[brazil]{babel}
\usepackage{microtype}

\usepackage[top=2.5cm, bottom=2.5cm, left=3cm, right=3cm, headheight=15pt]{geometry}

\usepackage{amsmath, amssymb, amsthm}
\usepackage{mathtools}
\usepackage{bm}

\usepackage{booktabs}
\usepackage{array}
\usepackage{multicol}
\usepackage{xcolor}
\usepackage{hyperref}
\usepackage{enumitem}
\usepackage{tcolorbox}
\tcbuselibrary{breakable}
\usepackage{fancyhdr}
\usepackage{titlesec}

% --------------------------------------------------------------------------
% Cores
% --------------------------------------------------------------------------
\definecolor{cyanrbc}{HTML}{3b9dbd}
\definecolor{orangerbc}{HTML}{d97b39}
\definecolor{grayrbc}{HTML}{6b7280}
\definecolor{greenbox}{HTML}{166534}
\definecolor{greenbg}{HTML}{dcfce7}
\definecolor{boxbg}{HTML}{f5f0e8}
\definecolor{boxborder}{HTML}{3b9dbd}
\definecolor{provabg}{HTML}{fff7ed}
\definecolor{provaborder}{HTML}{d97b39}
\definecolor{correctbg}{HTML}{dcfce7}
\definecolor{correctborder}{HTML}{166534}

% --------------------------------------------------------------------------
% Caixas
% --------------------------------------------------------------------------
\newtcolorbox{resultbox}[1][]{standard, colback=boxbg, colframe=boxborder,
  boxrule=1pt, arc=3pt, left=6pt, right=6pt, top=4pt, bottom=4pt,
  fonttitle=\bfseries\small, #1}
\newtcolorbox{provabox}[1][]{standard, colback=provabg, colframe=provaborder,
  boxrule=1.2pt, arc=3pt, left=6pt, right=6pt, top=4pt, bottom=4pt,
  fonttitle=\bfseries\small, #1}
\newtcolorbox{correctbox}[1][]{standard, colback=correctbg, colframe=correctborder,
  boxrule=1pt, arc=3pt, left=6pt, right=6pt, top=4pt, bottom=4pt,
  fonttitle=\bfseries\small\color{correctborder}, #1}
\newtcolorbox{formulabox}[1][]{standard, colback=greenbg!30, colframe=greenbox!60,
  boxrule=0.8pt, arc=3pt, left=6pt, right=6pt, top=4pt, bottom=4pt,
  fonttitle=\bfseries\small, #1}

% --------------------------------------------------------------------------
% Comandos matemáticos
% --------------------------------------------------------------------------
\newcommand{\E}{\mathbb{E}}
\newcommand{\dif}{\mathrm{d}}
\newcommand{\VAR}{\mathrm{VAR}}

\setlength{\parskip}{0.5\baselineskip}
\setlength{\parindent}{0pt}

\titleformat{\section}{\large\bfseries\color{cyanrbc}}{Questão \thesection.}{0.5em}{}[\vspace{-0.3em}\rule{\linewidth}{0.4pt}]
\titleformat{\subsection}{\normalsize\bfseries\color{grayrbc}}{}{0em}{}

% Cabeçalho de questão na parte "prova limpa"
\newcommand{\provaquestao}[2]{%
  \par\vspace{0.9em}%
  {\large\bfseries\color{cyanrbc} Questão #1.}\ {\bfseries #2}\par
  \vspace{-0.45em}\rule{\linewidth}{0.4pt}\par\vspace{0.1em}}

\pagestyle{fancy}
\fancyhf{}
\lhead{\footnotesize\textcolor{grayrbc}{Simulado P2 2026 --- Macroeconomia I (PPGEco-UFES)}}
\rhead{\footnotesize\textcolor{grayrbc}{Romer caps.\ 8 $\cdot$ 9 $\cdot$ 11 $\cdot$ 12}}
\cfoot{\small\thepage}
\renewcommand{\headrulewidth}{0.4pt}

% ==========================================================================
\begin{document}
% ==========================================================================

% --------------------------------------------------------------------------
% Capa
% --------------------------------------------------------------------------
\begin{titlepage}
\begin{center}
\vspace*{1.2cm}
{\Large\bfseries Universidade Federal do Espírito Santo}\\[0.5em]
{\large Programa de Pós-Graduação em Economia (PPGEco-UFES)}\\[2em]

\rule{\linewidth}{1.2pt}\\[0.8em]
{\LARGE\bfseries\color{cyanrbc} Simulado --- Prova 2\\[0.4em]
\large Macroeconomia I --- 2026}\\[0.4em]
{\normalsize\color{grayrbc} Romer caps.\ 8, 9, 11 e 12}\\[0.8em]
\rule{\linewidth}{1.2pt}

\vspace{1em}
{\large Modelo da P1: julgar as afirmativas e marcar a única correta}

\vspace{1.5em}
\begin{tcolorbox}[standard, colback=provabg, colframe=provaborder, width=0.9\textwidth,
  arc=4pt, boxrule=1.2pt, left=8pt, right=8pt, top=6pt, bottom=6pt]
\centering
\textbf{\textcolor{provaborder}{Conteúdo:}} 14 questões de escolha única (4 afirmativas; 1 correta)\\[0.4em]
\textbf{Consumo} (Q1--Q4) $\cdot$ \textbf{Investimento / $q$ de Tobin} (Q5--Q8)\\[0.2em]
\textbf{Inflação e Política Monetária} (Q9--Q11) $\cdot$
\textbf{Déficits e Política Fiscal} (Q12--Q14)
\end{tcolorbox}

\vspace{1.2em}
\begin{tcolorbox}[standard, colback=boxbg, colframe=boxborder, width=0.9\textwidth,
  arc=4pt, boxrule=1pt, left=8pt, right=8pt, top=6pt, bottom=6pt]
\raggedright
\textbf{Como usar este simulado (meta: $\geq 6{,}3$ na P2):}
\begin{enumerate}[leftmargin=1.5em, itemsep=0.15em, topsep=0.2em]
\item Resolva a \textbf{Parte I} sem consultar nada, cronometrando ($\approx$ 7 min/questão);
\item Confira no \textbf{gabarito rápido} e leia o comentário de \emph{todas} as
questões --- inclusive as que acertou: as alternativas falsas daqui são as
``pegadinhas'' que o professor reutiliza de prova em prova;
\item Cada questão indica, no rótulo \textbf{Radar}, de onde o tema saiu
(P2 2021, P2 2024, Lista 3 2026, slides). Questões recicladas nas provas
anteriores têm alta chance de reaparecer.
\end{enumerate}
\end{tcolorbox}

\vfill
{\small Simulado elaborado a partir de: P2 2021 (prova e respostas), P2 2024
(chave oficial de respostas), Lista 3 2026 (exercícios e chave oficial) e slides
das Aulas 10--13 (Romer, caps.\ 8, 9, 11 e 12).}
\end{center}
\end{titlepage}

% ==========================================================================
% PARTE I — A PROVA (sem respostas)
% ==========================================================================
\begin{center}
{\LARGE\bfseries\color{orangerbc} Parte I --- Prova}\\[0.3em]
{\normalsize\color{grayrbc} Julgue as afirmativas de cada questão e assinale a \textbf{única correta}.}\\[0.2em]
\rule{\linewidth}{0.8pt}
\end{center}

\begin{center}
{\large\bfseries\color{cyanrbc} Bloco I --- Teoria do Consumo (Cap.\ 8)}
\end{center}

% ------------------------------------------------------------------
\provaquestao{1}{Hipótese da Renda Permanente (Friedman, 1957)}

Na teoria do consumo de Friedman, a otimização da utilidade
$\sum_{t=1}^{T} u(C_t)$ sujeita à restrição orçamentária
$\sum_{t=1}^{T} C_t \leq A_0 + \sum_{t=1}^{T} Y_t$ (com $A_0$ a riqueza inicial
e taxa de juros nula) leva à solução
\[
C_t = \frac{1}{T}\Bigl(A_0 + \sum_{\tau=1}^{T} Y_\tau\Bigr), \qquad \text{para todo } t.
\]
Julgue os itens, escolhendo o correto:

\begin{enumerate}[label=\alph*)]
\item Sob a hipótese da renda permanente, o padrão temporal da renda ao longo
da vida é crítico para a determinação do consumo corrente, embora seja
irrelevante para a determinação da poupança.
\item Um aumento \emph{permanente} da renda (isto é, replicado em todos os
períodos até $T$) possui o mesmo impacto marginal $1/T$ sobre o consumo
corrente que um ganho puramente transitório.
\item Um incremento apenas corrente (transitório) da renda disponível possui
impacto marginal $1/T$ sobre o consumo corrente --- efeito decrescente no
horizonte de planejamento $T$ ---, o que implica propensão marginal a consumir
de curto prazo menor do que a prevista pela teoria keynesiana.
\item Como o consumo é determinado pela renda disponível corrente, regressões
do consumo sobre a renda tendem a estimar propensão marginal a consumir
próxima de 1 em amostras curtas.
\end{enumerate}

% ------------------------------------------------------------------
\provaquestao{2}{Passeio Aleatório (Hall, 1978) e o Teste de Campbell e Mankiw (1989)}

Considere a hipótese do passeio aleatório do consumo e a regressão proposta em
Campbell e Mankiw (1989),
\[
C_t - C_{t-1} = \lambda\,(Y_t - Y_{t-1}) + (1-\lambda)\,e_t.
\]
Julgue os itens, escolhendo o correto:

\begin{enumerate}[label=\alph*)]
\item Para Hall (1978), o consumo segue um passeio aleatório porque as famílias
ignoram a informação nova disponível em cada período ao formar suas
expectativas de consumo.
\item No teste com variáveis instrumentais de Campbell e Mankiw (dados
trimestrais dos EUA, 1953--1986), $\lambda$ mostrou-se estatisticamente
significativo, porém bem abaixo de 1 --- evidência de uma explicação
\emph{híbrida}: uma fração $\lambda$ das famílias consome segundo a renda
corrente (regra de bolso keynesiana) e a fração $(1-\lambda)$ comporta-se
conforme a renda permanente.
\item Campbell e Mankiw estimaram a regressão por Mínimos Quadrados Ordinários,
já que não há problema de endogeneidade na variável explicativa
$(Y_t - Y_{t-1})$.
\item O teste confirmou a rejeição plena da hipótese do passeio aleatório, com
estimativas convergindo para $\lambda = 1$.
\end{enumerate}

% ------------------------------------------------------------------
\provaquestao{3}{Equação de Euler: Juros Reais e Crescimento do Consumo}

Seja um modelo de otimização do consumo com taxa real de juros $r$, taxa de
desconto subjetiva $\rho$ e grau de aversão relativa ao risco $\theta$, cuja
equação de Euler é
\[
\frac{C_{t+1}}{C_t} = \left(\frac{1+r}{1+\rho}\right)^{\frac{1}{\theta}}.
\]
Julgue os itens, escolhendo o correto:

\begin{enumerate}[label=\alph*)]
\item Se $r < \rho$, o consumo tende a crescer ao longo do tempo.
\item A elevação do grau de aversão relativa ao risco aumenta a sensibilidade
do crescimento do consumo a variações da taxa real de juros.
\item O fato estilizado reportado pela literatura empírica (Hansen e Singleton,
1983; Hall, 1988b; Campbell e Mankiw, 1989) é claro: há grande impacto das
taxas reais de juros no comportamento do consumo das famílias.
\item O baixo poder explicativo das taxas reais de juros sobre o crescimento do
consumo, encontrado empiricamente para os EUA, é atribuído por esta literatura
a um grau elevado de aversão ao risco ($\theta$ alto) --- ou seja, a uma baixa
elasticidade de substituição intertemporal.
\end{enumerate}

% ------------------------------------------------------------------
\provaquestao{4}{Poupança Precaucional e Restrições de Liquidez}

A teoria moderna do consumo relaciona incerteza e poupança por meio de
\[
\E[g^c] \simeq \tfrac{1}{2}\,(\theta + 1)\,\VAR(g^c),
\]
sendo $\E[g^c]$ a expectativa de crescimento do consumo e $\VAR(g^c)$ a
variância dessa taxa de crescimento. Julgue os itens, escolhendo o correto:

\begin{enumerate}[label=\alph*)]
\item O aumento da incerteza quanto ao futuro ($\VAR(g^c)$ maior) reduz a
expectativa de crescimento do consumo, pois as famílias antecipam consumo para
o presente.
\item O \emph{puzzle} é que as famílias poupam relativamente pouco frente ao
que o motivo precaucional prevê; Carroll (1992, 1997) concilia teoria e
evidência apontando uma elevada taxa de desconto subjetiva, que atua reduzindo
a poupança e compensando o motivo precaucional.
\item Na expressão acima, o termo $(\theta + 1)$ representa o coeficiente de
aversão relativa ao risco da família.
\item Restrições de liquidez somente reduzem o consumo corrente quando estão
efetivamente ativas (\emph{binding}) no período presente.
\end{enumerate}

\begin{center}
{\large\bfseries\color{cyanrbc} Bloco II --- Teoria do Investimento e $q$ de Tobin (Cap.\ 9)}
\end{center}

% ------------------------------------------------------------------
\provaquestao{5}{Custos de Ajuste e a Condição de Otimalidade da Firma}

Considere o modelo de investimento com custos de ajuste do estoque de capital,
em que a firma maximiza
\[
\Pi = \int_{t=0}^{\infty} e^{-rt}\bigl[\pi(K(t))\,k(t) - I(t) - C(I(t))\bigr]\dif t.
\]
Julgue os itens, escolhendo o correto:

\begin{enumerate}[label=\alph*)]
\item Com custos de ajuste convexos --- $C(0)=0$, $C'(0)=0$ e $C''(\cdot)>0$ ---,
a maximização de lucros implica $1 + C'(I_t) = q_t$: a firma investe até o
ponto em que o custo de aquisição do capital (preço de compra mais custo
marginal de ajuste) iguala o valor do capital.
\item A condição de maximização de lucros implica que a firma investe até o
ponto em que o valor do capital supera o custo de sua aquisição.
\item O \emph{locus} $\dot{K}=0$ é caracterizado por $q = 0$.
\item Os custos de ajuste são funções lineares da taxa de investimento, o que
garante a existência de solução interior para o problema da firma.
\end{enumerate}

% ------------------------------------------------------------------
\provaquestao{6}{$q$ de Tobin: Determinação e Estática Comparativa}

O valor de mercado do bem de capital pode ser expresso por
\[
q(t) = \int_{t=0}^{\infty} e^{-rt}\,\pi(K(t))\,\dif t,
\qquad\text{com } \pi'(K) < 0,
\]
e, no \emph{locus} $\dot{q}=0$, vale $q = \pi(K)/r$. Julgue os itens,
escolhendo o correto:

\begin{enumerate}[label=\alph*)]
\item A expressão implica que o valor de mercado das empresas está
positivamente relacionado ao estoque de capital da indústria.
\item Uma expectativa de mercado de que a inflação deve permanecer acima da
meta nos próximos anos é condizente com maior $q$ da empresa, coeteris paribus,
pois isto é acompanhado de expectativa de maior crescimento do PIB e de maiores
fluxos futuros de caixa.
\item À medida que o estoque de capital $K$ da indústria aumenta, maior é o
valor $q$, de modo que a função $\dot{q}=0$ é positivamente inclinada no plano
$q \times K$.
\item Uma redução permanente da taxa real de juros eleva o $q$ associado a cada
nível de $K$ --- tornando a função $q = (1/r)\,\pi(K)$ mais inclinada --- e
desloca o equilíbrio de longo prazo da indústria para um estoque de capital
maior.
\end{enumerate}

% ------------------------------------------------------------------
\provaquestao{7}{Diagrama de Fases: $q$ e o Estoque de Capital}

Sejam as equações de movimento da indústria
\[
\dot{K}(t) = f\bigl(q(t)\bigr),\ \ f(1)=0,\ f'(\cdot)>0,
\qquad
\dot{q}(t) = r\,q(t) - \pi\bigl(K(t)\bigr).
\]
Julgue os itens, escolhendo o correto:

\begin{enumerate}[label=\alph*)]
\item O \emph{locus} abaixo de $\dot{K}=0$ e à direita de $\dot{q}=0$ é uma
região de divergência crescente em relação ao equilíbrio de longo prazo da
indústria.
\item Uma situação de estoque de capital muito baixo, acompanhada de nível
também muito reduzido para o $q$, implica elevação subsequente do $q$, dado que
os lucros marginais da indústria são altos.
\item Em uma situação acima do \emph{locus} $\dot{K}=0$ e à esquerda do
\emph{locus} $\dot{q}=0$, tem-se $q>1$ --- logo $\dot{K}>0$ --- e um $q$ abaixo
do valor consistente com os lucros correntes (pessimismo quanto ao futuro) ---
logo $\dot{q}<0$; a resultante dessas forças conduz a indústria, pela
trajetória de sela, ao equilíbrio de longo prazo $E$.
\item Sobre o \emph{locus} $\dot{K}=0$, o estoque de capital da indústria
cresce à taxa $r$.
\end{enumerate}

% ------------------------------------------------------------------
\provaquestao{8}{Choques de Produto, Juros e o Acelerador}

Quanto aos efeitos de variações do produto e das taxas de juros sobre o
investimento no modelo do $q$ de Tobin, julgue os itens, escolhendo o correto:

\begin{enumerate}[label=\alph*)]
\item Após uma elevação \emph{permanente} do produto, como o estoque de capital
não se ajusta instantaneamente, o $q$ salta e o investimento se eleva; à medida
que $K$ cresce, o preço relativo do produto da indústria e os lucros declinam,
até que o valor do capital retorne ao normal --- quando cessam os incentivos a
investimento adicional.
\item Enquanto uma redução permanente do produto vem acompanhada de diminuição
do investimento, dada a queda do $q$ de Tobin, uma redução transitória do
produto possui efeito nulo sobre o investimento.
\item O modelo do $q$ sustenta a visão de que as taxas de juros relevantes para
a decisão de investimento são as de curtíssimo prazo, e não as longas.
\item Pelo mecanismo do acelerador, o investimento é maior quando o produto
está em nível elevado há bastante tempo do que quando acabou de subir.
\end{enumerate}

\begin{center}
{\large\bfseries\color{cyanrbc} Bloco III --- Inflação e Política Monetária (Cap.\ 11)}
\end{center}

% ------------------------------------------------------------------
\provaquestao{9}{Mercado Monetário, Nível de Preços e Identidade de Fisher}

Considere o equilíbrio do mercado monetário e a identidade de Fisher,
\[
\frac{M}{P} = L(i, Y), \quad L_i < 0,\ L_Y > 0,
\qquad\qquad
i = r + \pi^e.
\]
Julgue os itens, escolhendo o correto:

\begin{enumerate}[label=\alph*)]
\item Grosso modo, uma alta de preços vem como consequência de variações
positivas na moeda, da taxa de juros e do produto real.
\item A identidade de Fisher estabelece que a taxa real de juros é a soma da
taxa nominal com a inflação esperada, $r = i + \pi^e$.
\item Sob plena flexibilidade de preços, variações na taxa de crescimento da
moeda alteram permanentemente o produto real e a taxa real de juros.
\item De $P = M/L(i,Y)$, o nível de preços eleva-se com aumentos da oferta de
moeda e da taxa de juros e com reduções do produto; não obstante, para a
inflação de longo prazo a ênfase recai sobre o crescimento da oferta monetária,
único fator capaz de gerar aumentos \emph{persistentes} do nível de preços.
\end{enumerate}

% ------------------------------------------------------------------
\provaquestao{10}{Crescimento da Moeda, Efeito Liquidez e Estrutura a Termo}

Sobre os efeitos de mudanças na taxa de crescimento da oferta monetária e a
estrutura a termo das taxas de juros, julgue os itens, escolhendo o correto:

\begin{enumerate}[label=\alph*)]
\item O efeito liquidez descreve a elevação imediata das taxas nominais de
curto prazo que se segue a uma expansão monetária.
\item Seja um modelo com plena flexibilidade de preços. Uma diminuição
permanente da taxa de crescimento da moeda pelo Banco Central implica queda
descontínua do nível de preços, da inflação esperada e da taxa nominal de
juros, de modo que haverá aumento do estoque real de moeda no longo prazo.
\item A estrutura a termo descreve a relação entre as taxas de juros para
distintos horizontes. O processo de arbitragem dos investidores satisfaz a
condição de que a taxa $i^n$ de um título de $n$ anos seja igual a $n\,i^1$,
sendo $i^1$ a taxa do título de 1 ano.
\item Pela teoria das expectativas da estrutura a termo, mudanças nas taxas
longas são determinadas por mudanças no prêmio a termo, e não por mudanças nas
expectativas das taxas curtas futuras.
\end{enumerate}

% ------------------------------------------------------------------
\provaquestao{11}{A Regra \emph{Forward-Looking} de Clarida, Galí e Gertler (2000)}

Considere a regra de política monetária
\[
i_t = r^n + \E_t[\pi_{t+n}]
      + \phi_\pi\bigl(\E_t[\pi_{t+n}] - \pi^*\bigr)
      + \phi_y\,\E_t\bigl[\ln Y_{t+n} - \ln Y^n_{t+n}\bigr],
\]
em que $r^n$ é a taxa real natural, $\pi^*$ a meta de inflação, e
$\ln Y_{t+n} - \ln Y^n_{t+n}$ o hiato esperado do produto. Julgue os itens,
escolhendo o correto:

\begin{enumerate}[label=\alph*)]
\item A regra é \emph{backward-looking}: o Banco Central reage aos desvios da
inflação passada em relação à meta, tal como na regra original de Taylor
(1993).
\item No equilíbrio de médio prazo --- quando $\E_t[\pi_{t+n}] = \pi^*$ e o
hiato esperado do produto é nulo ---, a regra implica $i_t = r^n - \pi^*$.
\item Sob o Princípio de Taylor ($\phi_\pi > 1$), a taxa nominal reage mais do
que proporcionalmente aos desvios da inflação esperada frente à meta, elevando
a taxa \emph{real} de juros após um choque inflacionário --- uma política
contracíclica que ancora expectativas e contribui para inflação baixa e
estável, a despeito do fato estilizado de correlação positiva de longo prazo
entre crescimento monetário e inflação.
\item Com $0 < \phi_\pi < 1$ a inflação também converge à meta, apenas mais
lentamente, pois qualquer resposta positiva da taxa nominal é suficiente para
elevar a taxa real de juros.
\end{enumerate}

\begin{center}
{\large\bfseries\color{cyanrbc} Bloco IV --- Déficits Orçamentários e Política Fiscal (Cap.\ 12)}
\end{center}

% ------------------------------------------------------------------
\provaquestao{12}{Restrição Orçamentária Intertemporal do Governo e Solvência}

A restrição orçamentária intertemporal do governo pode ser escrita como
\[
\int_{t=0}^{\infty} e^{-R(t)}\bigl[T(t) - G(t)\bigr]\dif t \;\geq\; D(0).
\]
Julgue os itens, escolhendo o correto:

\begin{enumerate}[label=\alph*)]
\item A avaliação quanto à sustentabilidade da dívida pública (sua solvência) é
feita a partir dos resultados fiscais correntes.
\item Dois países A e B possuem dívida pública real de mesma magnitude; no
entanto, os saldos fiscais correntes de A são muito superiores aos de B
($T_t^A - G_t^A > T_t^B - G_t^B$). Logo, a dívida pública de B será
necessariamente considerada mais arriscada que a de A.
\item Os superávits primários corrente e futuros esperados, a valor presente,
precisam cobrir a dívida pública atual em termos reais; por isso, anúncios e
sinalizações das autoridades fiscais têm especial importância na avaliação de
solvência, já que influenciam as projeções dos fluxos futuros de saldos
primários.
\item A satisfação da restrição exige que o resultado primário do governo seja
superavitário em todos os períodos do tempo.
\end{enumerate}

% ------------------------------------------------------------------
\provaquestao{13}{Equivalência Ricardiana}

Considere a restrição orçamentária das famílias no modelo RCK com tributos
\emph{lump-sum} e dívida pública,
\[
\int_{t=0}^{\infty} e^{-R(t)} C(t)\,\dif t \;\leq\;
K(0) + D(0) + \int_{t=0}^{\infty} e^{-R(t)}\bigl[W(t) - T(t)\bigr]\dif t,
\]
e a restrição do governo satisfeita com igualdade. Julgue os itens, escolhendo
o correto:

\begin{enumerate}[label=\alph*)]
\item Uma política de redução de alíquotas tributárias, ao elevar a renda
disponível corrente das famílias, vem acompanhada de aumento do consumo e,
desta maneira, de um estímulo real positivo para o ciclo econômico.
\item São os tributos, e não os gastos públicos, que afetam a acumulação de
capital, dado que $I(t) = Y(t) - C(t) - G(t)$.
\item A equivalência ricardiana permanece válida mesmo quando o governo não
satisfaz sua restrição orçamentária intertemporal.
\item Substituindo a restrição do governo na restrição das famílias, o consumo
passa a depender do valor presente dos \emph{gastos públicos} --- sem qualquer
referência à divisão de seu financiamento entre tributos e dívida; a renda
disponível extra gerada por um corte de tributos hoje, dados os gastos, é
poupada por agentes racionais que anteveem maiores tributos e superávits
futuros.
\end{enumerate}

% ------------------------------------------------------------------
\provaquestao{14}{Contrações Fiscais Expansionistas e Crises de Dívida}

Sobre a literatura empírica de política fiscal e os modelos de crises de
dívida, julgue os itens, escolhendo o correto:

\begin{enumerate}[label=\alph*)]
\item Alesina e Perotti (1997) mostram que consolidações fiscais baseadas em
aumentos de tributos são mais duradouras e mais frequentemente expansionistas
do que as baseadas em cortes de gastos com emprego público e transferências.
\item Giavazzi e Pagano (1990) documentam \emph{booms} de consumo após pacotes
de reforma fiscal na Dinamarca e na Irlanda nos anos 1980, atribuindo-os a
efeitos operando via expectativas --- canais: queda dos juros corrente e
esperados, menores distorções tributárias futuras (lado da oferta) e menor
probabilidade percebida de crise, elevando a renda futura esperada e, com ela,
consumo corrente e investimento (via $q$ de Tobin).
\item Nos modelos de crises de dívida inspirados em Calvo (1988) e Cole e Kehoe
(2000), uma dívida pública elevada reduz a taxa de retorno requerida pelos
credores, afastando a probabilidade de \emph{default}.
\item Consolidações fiscais são sempre contracionistas no curto prazo, uma vez
que os multiplicadores fiscais independem da composição do ajuste (gastos
\emph{versus} tributos).
\end{enumerate}

% ==========================================================================
% PARTE II — GABARITO
% ==========================================================================
\newpage
\begin{center}
{\LARGE\bfseries\color{orangerbc} Parte II --- Gabarito Comentado}\\[0.3em]
\rule{\linewidth}{0.8pt}
\end{center}

\begin{tcolorbox}[standard, colback=correctbg, colframe=correctborder,
  arc=4pt, boxrule=1pt, left=8pt, right=8pt, top=6pt, bottom=6pt]
\centering
\textbf{\textcolor{correctborder}{Gabarito Rápido}}\\[0.6em]
\begin{tabular}{ccccccc}
\toprule
Q1 & Q2 & Q3 & Q4 & Q5 & Q6 & Q7 \\
\textbf{c} & \textbf{b} & \textbf{d} & \textbf{b} & \textbf{a} & \textbf{d} & \textbf{c} \\
\midrule
Q8 & Q9 & Q10 & Q11 & Q12 & Q13 & Q14 \\
\textbf{a} & \textbf{d} & \textbf{b} & \textbf{c} & \textbf{c} & \textbf{d} & \textbf{b} \\
\bottomrule
\end{tabular}
\end{tcolorbox}

\begin{resultbox}[title={Gabaritos oficiais das provas anteriores (para referência)}]
\textbf{P2 2021} (itens corretos): Q1 = ii $\cdot$ Q2 = iv $\cdot$ Q3 = ii.\\
\textbf{P2 2024} (chave oficial): Q1 = ii $\cdot$ Q2 = iv $\cdot$ Q3 = iv $\cdot$
Q4 = iii $\cdot$ Q5 = iii $\cdot$ Q6 = iii.\\[0.3em]
O professor \emph{reutiliza} questões: a Q1 e a Q3 da P2 2024 são as mesmas
Q1 e Q2 da P2 2021 (com ajustes de redação). As alternativas falsas também se
repetem --- estude-as como estudaria as verdadeiras.
\end{resultbox}

% ==========================================================================
\newpage
\begin{center}
{\large\bfseries\color{cyanrbc} Bloco I --- Teoria do Consumo (Cap.\ 8)}\\[0.2em]
\rule{\linewidth}{0.6pt}
\end{center}

\setcounter{section}{0}

% ------------------------------------------------------------------
\section{Hipótese da Renda Permanente}
% ------------------------------------------------------------------

\subsection{Radar: Lista 3 2026, Q1 (chave oficial); P2 2021/2024, Q1; slides da Aula 10.}

\begin{correctbox}[title={Resposta correta: c)}]
Da condição de primeira ordem do Lagrangiano resulta
$C_t = \frac{1}{T}\bigl(A_0 + \sum_{\tau=1}^{T} Y_\tau\bigr)$: o consumo de
cada período é uma fração $1/T$ dos \emph{recursos de vida inteira}
(suavização).

Um aumento de $Y_1$ que não se replica nos demais períodos (componente
\textbf{transitório}) eleva o somatório em apenas uma parcela, com impacto
marginal $\partial C_t/\partial Y_1 = 1/T$ --- decrescente no horizonte $T$.
É exatamente o que diz a chave oficial da Lista 3: \emph{``um aumento da renda
disponível corrente $Y_1$ implica impacto marginal $1/T$ sobre o consumo
corrente. Este efeito é decrescente no horizonte de planejamento $T$''}. Por
isso a propensão marginal a consumir de curto prazo é menor que a keynesiana.
\checkmark

\textbf{Por que as demais são falsas:}
\begin{itemize}[leftmargin=*]
\item[a)] Inverte o resultado dos slides: o padrão temporal da renda
\emph{não} é importante para o consumo (que depende só do valor total dos
recursos), mas é \emph{crítico} para a poupança, pois
$S_t = Y_t - C_t = \bigl(Y_t - \frac{1}{T}\sum_\tau Y_\tau\bigr) - \frac{1}{T}A_0$.
\item[b)] Um aumento permanente (em todos os $T$ períodos) soma $T$ parcelas
iguais: impacto total $\approx T \cdot (1/T) = 1$, ou seja, o consumo sobe
praticamente um para um --- e não $1/T$.
\item[d)] Em amostras \emph{curtas}, a variância da renda transitória domina e
o coeficiente estimado
$\hat{b} = \frac{\VAR(Y^P)}{\VAR(Y^P) + \VAR(Y^T)}$
fica bem \emph{abaixo} de 1; é em amostras longas que a variância permanente
domina e $\hat{b} \to 1$.
\end{itemize}
\end{correctbox}

\begin{formulabox}[title={Para memorizar}]
$C_t = \dfrac{1}{T}\Bigl(A_0 + \sum_{\tau=1}^{T} Y_\tau\Bigr)$
\quad$\cdot$\quad transitório: impacto $1/T$
\quad$\cdot$\quad permanente: impacto $\approx 1$
\quad$\cdot$\quad $\hat{b} = \dfrac{\VAR(Y^P)}{\VAR(Y^P)+\VAR(Y^T)}$
\end{formulabox}

% ------------------------------------------------------------------
\section{Hall (1978) e Campbell e Mankiw (1989)}
% ------------------------------------------------------------------

\subsection{Radar: item correto da P2 2021/2024 Q1 (= ii); chave da Lista 3 2026, Q1b--c.}

\begin{correctbox}[title={Resposta correta: b)}]
Campbell e Mankiw (1989) testam Hall contra a alternativa de que uma fração
$\lambda$ dos consumidores gasta a renda corrente e o restante segue a renda
permanente. Estimando por \textbf{variáveis instrumentais} (para controlar a
endogeneidade de $\Delta Y_t$), com dados trimestrais de 1953--1986, obtêm
$\lambda$ \textbf{estatisticamente significativo, porém bem abaixo de 1}
($\approx 0{,}5$). Interpretação da chave oficial: \emph{explicação híbrida}
--- parte das famílias à la Keynes, parte $(1-\lambda)$ à la renda
permanente/passeio aleatório. \checkmark

\textbf{Por que as demais são falsas:}
\begin{itemize}[leftmargin=*]
\item[a)] É o oposto: o consumo segue passeio aleatório justamente porque
\emph{toda a informação disponível já está incorporada} ao consumo corrente
($C_t = \E_t[C_{t+1}]$, logo $C_t = C_{t-1} + e_t$); variações futuras só
decorrem de informação \emph{nova} sobre a renda permanente.
\item[c)] Há endogeneidade em $\Delta Y_t$ (correlação com o erro), e é
precisamente por isso que os autores usam IV --- não MQO.
\item[d)] Armadilha clássica do professor (P2 2021 e 2024): a rejeição
\emph{plena} do passeio aleatório exigiria $\lambda = 1$; o teste encontrou
$\lambda$ bem abaixo de 1, ou seja, rejeição apenas \emph{parcial}.
\end{itemize}
\end{correctbox}

% ------------------------------------------------------------------
\section{Equação de Euler: Juros e Consumo}
% ------------------------------------------------------------------

\subsection{Radar: Lista 3 2026, Q2 (chave oficial); P2 2024, Q2; slides da Aula 10.}

\begin{correctbox}[title={Resposta correta: d)}]
Da equação de Euler, a elasticidade do crescimento do consumo à taxa real de
juros é $1/\theta$ (elasticidade de substituição intertemporal, EIS). A chave
oficial da Lista 3 é direta: \emph{``os autores mencionados associam o fato
empírico de baixo poder preditivo de juros reais sobre $C_{t+1}/C_t$ ao suposto
grau elevado de aversão ao risco ($\theta$)''}. $\theta$ alto $\Rightarrow$
EIS $= 1/\theta$ baixa $\Rightarrow$ consumo pouco sensível a $r$. \checkmark

\textbf{Por que as demais são falsas:}
\begin{itemize}[leftmargin=*]
\item[a)] $r < \rho \Rightarrow (1+r)/(1+\rho) < 1 \Rightarrow C_{t+1}/C_t < 1$:
o consumo tende a \emph{declinar} (impaciência domina o prêmio por poupar).
\item[b)] Sensibilidade é $1/\theta$: aversão maior \emph{reduz} a resposta do
consumo aos juros.
\item[c)] O fato estilizado é o \emph{baixo} impacto empírico dos juros reais
sobre o crescimento do consumo --- este item (falso) apareceu textualmente como
distrator na P2 2024, Q2-ii.
\end{itemize}
\end{correctbox}

% ------------------------------------------------------------------
\section{Poupança Precaucional e Restrições de Liquidez}
% ------------------------------------------------------------------

\subsection{Radar: Lista 3 2026, Q3 (chave oficial); distrator recorrente na P2 2021/2024 Q1-iv; slides da Aula 10 (Zeldes; Gross e Souleles).}

\begin{correctbox}[title={Resposta correta: b)}]
A teoria prevê que mais incerteza ($\VAR(g^c)\uparrow$) induz poupança
precaucional e, com ela, maior crescimento esperado do consumo. O
\textbf{fato estilizado}, porém, é que \emph{``as famílias poupam relativamente
pouco, ou abaixo do que o modelo pressupõe''} (chave oficial). A conciliação é
de Carroll (1992, 1997): a \textbf{elevada taxa de desconto subjetiva} atua
diminuindo a poupança, \emph{``anulando o aumento da poupança pelo motivo
precaucional''}. \checkmark

\textbf{Por que as demais são falsas:}
\begin{itemize}[leftmargin=*]
\item[a)] Pela expressão, $\VAR(g^c)\uparrow \Rightarrow \E[g^c]\uparrow$: a
incerteza \emph{eleva} (não reduz) o crescimento esperado do consumo --- as
famílias adiam consumo (poupam mais) para se proteger.
\item[c)] $(\theta+1)$ é o coeficiente de \textbf{prudência relativa} (nos
slides, $\delta$), não o de aversão ao risco ($\theta$). São conceitos
distintos: prudência é o que gera o motivo precaucional ($u''' > 0$).
\item[d)] Zeldes (1989): mesmo restrições \emph{não} ativas hoje reduzem o
consumo corrente, se houver possibilidade de que venham a ser ativas no futuro
(Gross e Souleles, 2002). $\uparrow$restrição (efetiva ou esperada)
$\Rightarrow \downarrow$consumo.
\end{itemize}
\end{correctbox}

\begin{formulabox}[title={Atenção à armadilha da P2 2021/2024 (Q1-iv)}]
O item que associa poupança precaucional ao \emph{``fato estilizado de uma
elevada poupança pelas famílias''} é \textbf{falso}: o fato estilizado é de
poupança \emph{baixa} --- este é o puzzle que Carroll resolve.
\end{formulabox}

% ==========================================================================
\newpage
\begin{center}
{\large\bfseries\color{cyanrbc} Bloco II --- Investimento e $q$ de Tobin (Cap.\ 9)}\\[0.2em]
\rule{\linewidth}{0.6pt}
\end{center}

% ------------------------------------------------------------------
\section{Custos de Ajuste e Otimalidade da Firma}
% ------------------------------------------------------------------

\subsection{Radar: distrator textual da P2 2021 Q2-i e P2 2024 Q3-i; slides da Aula 11.}

\begin{correctbox}[title={Resposta correta: a)}]
Com custos de ajuste convexos, a condição de primeira ordem do problema da
firma é
\[
1 + C'(I_t) = q_t,
\]
lida nos slides como: \emph{``a firma investe até o ponto em que o custo de
aquisição do capital (preço de compra --- fixado em 1 --- mais o custo marginal
de ajuste) \textbf{iguala} o valor do capital ($q_t$)''}. \checkmark

\textbf{Por que as demais são falsas:}
\begin{itemize}[leftmargin=*]
\item[b)] Distrator usado nas duas provas anteriores: no ótimo há
\emph{igualdade}, não excesso do valor sobre o custo. Se o valor superasse o
custo, haveria lucro marginal não explorado --- a firma investiria mais.
\item[c)] $\dot{K} = f(q)$ com $f(1) = 0$: o \emph{locus} $\dot{K}=0$
corresponde a $q = 1$ (não $q=0$): vale investir se $q>1$ e desinvestir se
$q<1$.
\item[d)] Os custos de ajuste são \emph{convexos} ($C''>0$, com $C(0)=0$ e
$C'(0)=0$); é a convexidade que suaviza o investimento e gera solução interior.
Custos lineares levariam a ajuste instantâneo (solução de canto).
\end{itemize}
\end{correctbox}

% ------------------------------------------------------------------
\section{$q$ de Tobin: Determinação e Estática Comparativa}
% ------------------------------------------------------------------

\subsection{Radar: P2 2024, Q4 (gabarito oficial: iii); Lista 3 2026, Q4 (chave); slides da Aula 11 (Figura 9.7).}

\begin{correctbox}[title={Resposta correta: d)}]
No \emph{locus} $\dot{q}=0$, $q = \pi(K)/r$. Uma queda permanente de $r$:
(i) eleva o valor presente dos fluxos futuros de lucros --- $q$ maior para cada
$K$; (ii) torna a função $q = (1/r)\pi(K)$ \emph{mais inclinada} (slides,
Figura 9.7); (iii) desloca o equilíbrio $E$ para a direita: com $q$
momentaneamente acima de 1, $\dot{K} > 0$ até o novo equilíbrio com $K$ maior.
É também a resposta da chave da Lista 3, Q4b: juros esperados têm relação
\emph{inversa} com o $q$. \checkmark

\textbf{Por que as demais são falsas:}
\begin{itemize}[leftmargin=*]
\item[a)] $\pi'(K) < 0$: mais capital na indústria $\Rightarrow$ lucros
marginais menores $\Rightarrow$ valor de mercado \emph{menor}. Relação
\emph{negativa} (distrator da P2 2024, Q4-i).
\item[b)] Distrator da P2 2024, Q4-ii: inflação esperada \emph{acima da meta}
leva o Banco Central (regra com $\phi_\pi > 1$) a elevar os juros reais ---
maior taxa de desconto $\Rightarrow$ $q$ \emph{menor}, e não maior.
\item[c)] Como $\pi'(K)<0$, o \emph{locus} $\dot{q}=0$ é \emph{negativamente}
inclinado no plano $q \times K$ (distrator da P2 2024, Q4-iv).
\end{itemize}
\end{correctbox}

\begin{formulabox}[title={Para memorizar}]
$q(t) = \displaystyle\int_0^\infty e^{-rt}\pi(K(t))\,\dif t$
\quad$\cdot$\quad $\dot{q} = rq - \pi(K)$
\quad$\cdot$\quad $\dot{q}=0 \Rightarrow q = \pi(K)/r$ (inclinação negativa)
\quad$\cdot$\quad $\dot{K}=0 \Rightarrow q = 1$
\end{formulabox}

% ------------------------------------------------------------------
\section{Diagrama de Fases: $q$ e Capital}
% ------------------------------------------------------------------

\subsection{Radar: Lista 3 2026, Q5 (chave oficial); P2 2021 Q2 / P2 2024 Q3 (gabarito: iv); slides da Aula 11 (Figuras 9.1--9.4).}

\begin{correctbox}[title={Resposta correta: c)}]
É a situação da Q5 da Lista 3 2026, e a resposta reproduz a chave oficial:
acima de $\dot{K}=0$ tem-se $q > 1$, logo
$\dot{K} = f(q) > 0$ (capital crescendo); à esquerda de $\dot{q}=0$, o $q$ está
\emph{abaixo} do valor sugerido pelos lucros correntes
($q < \pi(K)/r$, ``pessimismo quanto ao futuro''), logo
$\dot{q} = rq - \pi(K) < 0$. \emph{``A resultante dessas forças cria um vetor
que leva a indústria rumo ao ponto E''} --- é a região da
\textbf{trajetória de sela}: convergência. \checkmark

\textbf{Por que as demais são falsas:}
\begin{itemize}[leftmargin=*]
\item[a)] Distrator textual das provas 2021/2024 (item ii, sempre falso): a
região abaixo de $\dot{K}=0$ e à direita de $\dot{q}=0$ ($q<1$, $K$ alto,
$\dot{K}<0$, $\dot{q}>0$) contém o \emph{braço convergente} da trajetória de
sela vindo do sudeste --- não é região de divergência crescente.
\item[b)] Inversão do item \emph{verdadeiro} das provas 2021/2024 (Q2-iv/Q3-iv):
com $K$ muito baixo e $q$ muito baixo, a equação $\dot{q} = rq - \pi(K)$ dá
$rq$ pequeno e $\pi(K)$ grande $\Rightarrow$ $\dot{q} < 0$: o $q$ sofre
\emph{queda adicional} (região de divergência abaixo da sela), e não elevação.
\item[d)] Sobre $\dot{K}=0$, por definição, o estoque de capital está
\emph{estável} ($q=1$, investimento líquido nulo) --- não cresce a taxa
alguma.
\end{itemize}
\end{correctbox}

% ------------------------------------------------------------------
\section{Choques de Produto, Juros e o Acelerador}
% ------------------------------------------------------------------

\subsection{Radar: slides da Aula 11 (Figuras 9.5 e 9.6); distrator da P2 2021 Q2-iii / P2 2024 Q3-iii; Lista 3 2026, Q4a.}

\begin{correctbox}[title={Resposta correta: a)}]
Narrativa exata dos slides (Figura 9.5, aumento permanente do produto): a
demanda pelo produto da indústria sobe; como $K$ não se ajusta
instantaneamente, o valor de mercado do capital ($q$) salta e atrai
investimento; com o crescimento de $K$, o preço relativo do produto e os lucros
vão declinando; o processo continua \emph{``até que o valor do capital retorne
ao normal, quando não há incentivos para investimento adicional''} (novo
equilíbrio $E'$ com $K$ maior e $q$ de volta a 1). \checkmark

\textbf{Por que as demais são falsas:}
\begin{itemize}[leftmargin=*]
\item[b)] Distrator recorrente (2021/2024): o choque \emph{transitório} tem
efeito \textbf{menor} --- o $q$ sobe menos porque os lucros extras duram pouco
---, mas \emph{não nulo}.
\item[c)] Os slides são explícitos: o modelo sustenta a visão de que as taxas
de juros \textbf{de longo prazo} são as relevantes para o investimento
($r$ na fórmula do $q$ é uma taxa longa).
\item[d)] Acelerador invertido. Correto: \emph{``o investimento é maior quando
o produto subiu recentemente do que quando está alto há um período
prolongado''} --- é a \emph{variação} do produto que move o investimento.
\end{itemize}
\end{correctbox}

% ==========================================================================
\newpage
\begin{center}
{\large\bfseries\color{cyanrbc} Bloco III --- Inflação e Política Monetária (Cap.\ 11)}\\[0.2em]
\rule{\linewidth}{0.6pt}
\end{center}

% ------------------------------------------------------------------
\section{Mercado Monetário, Preços e Fisher}
% ------------------------------------------------------------------

\subsection{Radar: P2 2021, Q3-i (distrator); slides da Aula 12 (equações 1--4).}

\begin{correctbox}[title={Resposta correta: d)}]
De $P = M/L(i,Y)$: $M\uparrow \Rightarrow P\uparrow$;
$i\uparrow \Rightarrow L\downarrow \Rightarrow P\uparrow$;
$Y\downarrow \Rightarrow L\downarrow \Rightarrow P\uparrow$. E, como nos
slides: \emph{``para entender a inflação no longo prazo, os economistas
enfatizam um único fator: o crescimento da oferta monetária''}, pois nenhum
outro fator gera aumentos \emph{persistentes} do nível de preços (os demais são
variações de nível, limitadas). \checkmark

\textbf{Por que as demais são falsas:}
\begin{itemize}[leftmargin=*]
\item[a)] Distrator da P2 2021 (Q3-i): o produto real entra
\emph{negativamente} --- $Y\uparrow$ eleva a demanda por moeda e \emph{reduz} o
nível de preços, dado $M$.
\item[b)] Identidade de Fisher: $i = r + \pi^e$ (ou $r = i - \pi^e$). O item
soma no lado errado.
\item[c)] Com preços plenamente flexíveis vale a \emph{neutralidade}: moeda não
afeta produto real nem taxa real de juros (constantes em $\bar{Y}$ e
$\bar{r}$); o crescimento monetário vai para inflação e juro nominal
(efeito Fisher).
\end{itemize}
\end{correctbox}

% ------------------------------------------------------------------
\section{Crescimento da Moeda, Efeito Liquidez e Estrutura a Termo}
% ------------------------------------------------------------------

\subsection{Radar: resposta correta da P2 2021, Q3 (= ii, aqui com sinal invertido: redução em vez de aumento); distratores da P2 2021 Q3-iii/iv; slides da Aula 12.}

\begin{correctbox}[title={Resposta correta: b)}]
É a resposta oficial da P2 2021 (Q3-ii). Com preços flexíveis, uma
\textbf{redução permanente do crescimento da moeda}:
$\Delta M\downarrow \Rightarrow \pi^e\downarrow \Rightarrow i\downarrow
\Rightarrow L(i,Y)\uparrow$: a demanda por moeda real \emph{sobe}
descontinuamente; como $M$ não salta, o \emph{nível de preços} cai
descontinuamente no anúncio, e no longo prazo o estoque real de moeda $M/P$ fica
\emph{maior}. (Espelho exato da Figura 11.2 dos slides, que mostra o caso de
\emph{aumento} do crescimento da moeda: $P$ salta para cima e $M/P$ cai.)
\checkmark

\textbf{Por que as demais são falsas:}
\begin{itemize}[leftmargin=*]
\item[a)] O efeito liquidez é a \emph{redução} imediata das taxas nominais
curtas após expansão monetária (com preços rígidos:
$\uparrow M/P \Rightarrow \downarrow r \Rightarrow \downarrow i
\Rightarrow \uparrow Y$).
\item[c)] Distrator da P2 2021 (Q3-iii): a arbitragem implica que $i^n$ é a
\emph{média} das taxas curtas esperadas mais o prêmio a termo,
\[
i_t^n = \frac{i_t^1 + \E_t i_{t+1}^1 + \cdots + \E_t i_{t+n-1}^1}{n} + \theta_{nt},
\]
e não $n\,i^1$.
\item[d)] A teoria das expectativas diz o \emph{contrário}: mudanças na
estrutura a termo refletem mudanças nas \textbf{expectativas de taxas curtas
futuras}, não no prêmio a termo. (Consequência dos slides: expansão monetária
tende a \emph{reduzir} as taxas curtas --- efeito liquidez --- e \emph{elevar}
as longas, que já embutem a inflação futura.)
\end{itemize}
\end{correctbox}

% ------------------------------------------------------------------
\section{Regra de Clarida, Galí e Gertler (2000)}
% ------------------------------------------------------------------

\subsection{Radar: novidade da Lista 3 2026 (Q6, com chave oficial detalhada) --- forte candidata à P2; slides da Aula 12 (calibração do BCB: $\phi_\pi = 1{,}5$; $\phi_y = 0{,}3$).}

\begin{correctbox}[title={Resposta correta: c)}]
É o \textbf{Princípio de Taylor}, coração da chave oficial da Q6 da Lista 3:
com $\phi_\pi > 1$, um aumento da inflação esperada eleva a taxa nominal
\emph{mais do que proporcionalmente}; como
$\Delta r_t = \Delta i_t - \Delta \E_t(\pi_{t+n}) > 0$, a taxa \emph{real} sobe
após o choque inflacionário, contraindo a demanda agregada e trazendo a
inflação de volta à meta --- política \textbf{contracíclica} que ancora
expectativas. É assim que economias que adotam a regra evitam o fato estilizado
de inflação média alta associada a crescimento monetário alto. \checkmark

\textbf{Por que as demais são falsas:}
\begin{itemize}[leftmargin=*]
\item[a)] A regra é \emph{forward-looking}: reage a $\E_t[\pi_{t+n}]$
(expectativas) e ao hiato \emph{esperado} --- é uma extensão da regra de Taylor
(1993) com visão prospectiva do Banco Central.
\item[b)] No equilíbrio de médio prazo os termos de desvio zeram e
$i_t = r^n + \E_t[\pi_{t+n}] = r^n + \pi^*$ (soma, não subtração). Com a
calibração dos slides: $i = 5{,}0 + 3{,}0 = 8{,}0\%$.
\item[d)] Com $\phi_\pi < 1$ a taxa nominal sobe \emph{menos} que a inflação
esperada: a taxa real \emph{cai} (ou não sobe o suficiente), a política
torna-se acomodatícia, a inflação fica persistente e as expectativas podem se
\textbf{desancorar} --- não há convergência ``apenas mais lenta''.
\end{itemize}
\end{correctbox}

\begin{formulabox}[title={Para memorizar (chave da Lista 3, Q6)}]
$i_t = r^n + \E_t[\pi_{t+n}] + \phi_\pi\bigl(\E_t[\pi_{t+n}] - \pi^*\bigr)
+ \phi_y\,\E_t[\ln Y_{t+n} - \ln Y^n_{t+n}]$\\[0.3em]
Princípio de Taylor: $\phi_\pi > 1 \Rightarrow$ juro real sobe com a inflação
esperada $\Rightarrow$ inflação baixa e estável.\\
Calibração citada em aula (BCB): $r^n = 5\%$, $\pi^* = 3\%$, $n = 4$ trimestres,
$\phi_\pi = 1{,}5$, $\phi_y = 0{,}3$.
\end{formulabox}

% ==========================================================================
\newpage
\begin{center}
{\large\bfseries\color{cyanrbc} Bloco IV --- Déficits e Política Fiscal (Cap.\ 12)}\\[0.2em]
\rule{\linewidth}{0.6pt}
\end{center}

% ------------------------------------------------------------------
\section{Restrição Orçamentária do Governo e Solvência}
% ------------------------------------------------------------------

\subsection{Radar: P2 2024, Q6 (gabarito oficial: iii); Lista 3 2026, Q7 (chave oficial).}

\begin{correctbox}[title={Resposta correta: c)}]
É o item verdadeiro da P2 2024 (Q6-iii), fundamentado na condição de solvência
da chave da Lista 3:
\[
\int_{t=0}^{\infty} e^{-R(t)}\bigl[T(t) - G(t)\bigr]\dif t \;\geq\; D(0).
\]
A solvência é avaliada pela capacidade de gerar fluxos de superávits primários
--- \emph{corrente e futuros esperados}, a valor presente --- suficientes para
cobrir a dívida atual. Como o resultado corrente é só uma pequena fração do
lado esquerdo, \emph{``anúncios de política fiscal e sinalizações por parte das
autoridades fiscais têm especial importância na avaliação de solvência''}
(expectativas, reputação e credibilidade). \checkmark

\textbf{Por que as demais são falsas:}
\begin{itemize}[leftmargin=*]
\item[a)] Distrator da P2 2024 (Q6-i): a sustentabilidade depende dos fluxos
\emph{projetados}, não apenas dos resultados correntes.
\item[b)] Distrator da P2 2024 (Q6-ii): o ``necessariamente'' derruba o item
--- B pode ter fluxos futuros esperados melhores (reformas anunciadas,
crescimento) e ser percebido como \emph{menos} arriscado apesar do corrente
pior.
\item[d)] A restrição vale para o \emph{somatório a valor presente}: déficits
primários temporários são perfeitamente compatíveis com solvência, desde que
compensados por superávits futuros esperados.
\end{itemize}
\end{correctbox}

% ------------------------------------------------------------------
\section{Equivalência Ricardiana}
% ------------------------------------------------------------------

\subsection{Radar: Lista 3 2026, Q8 (chave oficial, com derivação completa); P2 2024, Q5 (gabarito: iii) e Q2-iv; dissertativa recorrente (P2 2021 e 2024).}

\begin{correctbox}[title={Resposta correta: d)}]
Derivação da chave oficial (Lista 3, Q8): a restrição do governo satisfeita com
igualdade dá
$\int e^{-R}T\,\dif t = \int e^{-R}G\,\dif t + D(0)$;
substituindo na restrição das famílias,
\[
\int_{t=0}^{\infty} e^{-R(t)} C(t)\,\dif t \;\leq\;
K(0) + \int_{t=0}^{\infty} e^{-R(t)} W(t)\,\dif t
- \int_{t=0}^{\infty} e^{-R(t)} G(t)\,\dif t.
\]
$D(0)$ e a trajetória de $T$ \textbf{desaparecem}: só o valor presente dos
\emph{gastos} importa. Corte de tributos hoje, dados os gastos, eleva a dívida;
agentes racionais anteveem maiores tributos/superávits futuros e \emph{poupam}
a renda extra (transitória) --- efeito nulo sobre consumo e atividade.
\checkmark

\textbf{Por que as demais são falsas:}
\begin{itemize}[leftmargin=*]
\item[a)] Distrator da P2 2024 (Q5-ii): sob equivalência ricardiana o corte de
alíquotas \emph{não} estimula o consumo --- a renda extra é poupada.
\item[b)] Inversão da implicação 2 dos slides: \emph{``são os gastos públicos,
não os tributos, que afetam a acumulação de capital''}, pois
$I = Y - C - G$.
\item[c)] A derivação \emph{exige} a restrição do governo satisfeita (com
igualdade); sem ela não se pode substituir $\int e^{-R}T$ e o resultado não se
sustenta.
\end{itemize}
\end{correctbox}

% ------------------------------------------------------------------
\section{Contrações Fiscais Expansionistas e Crises de Dívida}
% ------------------------------------------------------------------

\subsection{Radar: slides da Aula 13 (Giavazzi-Pagano; Alesina-Perotti; Calvo; Cole-Kehoe) --- conteúdo novo, ainda não cobrado nas P2 anteriores.}

\begin{correctbox}[title={Resposta correta: b)}]
Giavazzi e Pagano (1990): os pacotes de reforma fiscal da \textbf{Dinamarca e
da Irlanda nos anos 1980} foram seguidos de \emph{booms de consumo}, atribuídos
a efeitos via \textbf{expectativas}. Canais dos slides: (a) juros --- menos
gasto público corrente e esperado reduz juros correntes e esperados; (b) lado
da oferta --- menores tributos futuros implicam menores distorções e maior
renda futura; (c) menor probabilidade percebida de crise; (d) maior renda
futura esperada eleva consumo corrente e investimento (Bertola e Drazen, 1993;
Perotti, 1999; via $q$ de Tobin). \checkmark

\textbf{Por que as demais são falsas:}
\begin{itemize}[leftmargin=*]
\item[a)] Inversão de Alesina e Perotti (1997): são as consolidações via
\emph{cortes de gastos} (emprego público e transferências) que tendem a ser
mais duradouras e frequentemente \emph{expansionistas}; as baseadas em aumentos
de tributos usualmente não o são.
\item[c)] Nos modelos de crises de dívida (Calvo, 1988; Cole e Kehoe, 2000) a
lógica é: \emph{dívida alta $\rightarrow$ taxa de retorno requerida alta
$\rightarrow$ receitas futuras esperadas baixas $\rightarrow$ default mais
provável} --- crises \textbf{induzidas por expectativas} (profecias
autorrealizáveis), o oposto do afirmado.
\item[d)] Falso duplamente: multiplicadores \emph{dependem} da composição do
ajuste, e consolidações podem ser expansionistas (é o ponto central da
literatura acima).
\end{itemize}
\end{correctbox}

% ==========================================================================
\newpage
\begin{center}
{\large\bfseries\color{orangerbc} Checklist final de véspera}\\[0.2em]
\rule{\linewidth}{0.6pt}
\end{center}

\begin{resultbox}[title={As 10 armadilhas que o professor mais usa (todas já caíram)}]
\begin{enumerate}[leftmargin=1.6em, itemsep=0.2em]
\item \textbf{$\lambda = 1$}: rejeição \emph{plena} do passeio aleatório
exigiria $\lambda=1$; Campbell-Mankiw acharam $\lambda$ significativo mas
$\ll 1$ (híbrido).
\item \textbf{Poupança precaucional}: o fato estilizado é poupança
\emph{baixa} (puzzle; Carroll e a alta taxa de desconto), nunca ``elevada
poupança''.
\item \textbf{Juros e consumo}: relação teórica \emph{positiva}, impacto
empírico \emph{fraco} ($\theta$ alto / EIS baixa).
\item \textbf{Firma investe até valor = custo} ($1 + C'(I) = q$), nunca
``supera''.
\item \textbf{Choque transitório de produto}: efeito \emph{menor}, jamais
\emph{nulo}.
\item \textbf{$q$ vs.\ $K$}: $\pi'(K)<0 \Rightarrow$ \emph{locus} $\dot q = 0$
com inclinação \emph{negativa}; valor da empresa cai com $K$ da indústria.
\item \textbf{$K$ baixo e $q$ baixo} $\Rightarrow \dot q = rq - \pi(K) < 0$:
$q$ cai ainda mais (região divergente).
\item \textbf{Preços}: $P = M/L(i,Y)$ --- produto entra \emph{negativo};
inflação persistente $=$ crescimento da moeda.
\item \textbf{Estrutura a termo}: $i^n$ é \emph{média} das curtas esperadas
$+$ prêmio; expansão monetária derruba as curtas (liquidez) e sobe as longas.
\item \textbf{Princípio de Taylor}: $\phi_\pi > 1$ para o juro \emph{real}
subir; $\phi_\pi<1 \Rightarrow$ persistência e desancoragem.
\end{enumerate}
\end{resultbox}

\begin{provabox}[title={Prognóstico para a P2 2026 (baseado no padrão 2021 $\to$ 2024 $\to$ lista 2026)}]
\begin{itemize}[leftmargin=1.5em, itemsep=0.15em]
\item \textbf{Certeza quase total}: 1 questão de HRP/Hall/Campbell-Mankiw
(Q1 reciclada em 2021 e 2024); 1 de investimento com diagrama de fases
(Q2/Q3 recicladas); 1 de equivalência ricardiana; 1 de restrição
orçamentária/solvência.
\item \textbf{Muito provável}: 1 questão sobre a regra de Clarida et al.\
(2000) --- é a \emph{novidade} da Lista 3 2026 (Q6) e ganhou chave oficial
detalhada; Princípio de Taylor no centro.
\item \textbf{Provável}: $q$ de Tobin conceitual (P2 2024 Q4); mercado
monetário/efeito liquidez/estrutura a termo (P2 2021 Q3); poupança
precaucional; contrações expansionistas (slides novos).
\item \textbf{Formato}: modelo P1 --- julgar afirmativas e marcar a única
correta. Leia \emph{todas} as alternativas: o professor troca uma palavra
(``supera''/``iguala''; ``nulo''/``menor''; ``positiva''/``negativa'') para
inverter o valor-verdade.
\end{itemize}
\end{provabox}

\vfill
{\small\color{grayrbc} Simulado de estudo --- Macroeconomia I, PPGEco-UFES 2026.
Elaborado a partir dos materiais oficiais da disciplina (P2 2021, P2 2024 com
chave, Lista 3 2026 com chave, slides das Aulas 10--13). Uso pessoal.}

\end{document}
